\chapter*{术语表}
\addcontentsline{toc}{chapter}{术语表}

\begin{description}
  \item[张量] 带元素类型、形状、步长和生命周期的数据视图。
  \item[算子] 有输入合同、输出合同和错误边界的计算单元。
  \item[CPU] Central Processing Unit，中央处理器。本册 CPU 后端通常指本地内存、SIMD、多线程和 cache/TLB 下的实现。
  \item[GPU] Graphics Processing Unit，图形处理器或通用并行加速器。本册 GPU 后端关注 kernel、stream、显存和同步边界。
  \item[ISA] Instruction Set Architecture，指令集架构，规定软件可见的指令、寄存器和状态。
  \item[API] Application Programming Interface，源码层调用接口。
  \item[ABI] Application Binary Interface，二进制层调用、对象布局和链接合同。
  \item[SDK] Software Development Kit，软件开发工具包，通常包含头文件、库、工具、示例和文档。
  \item[GDB] GNU Debugger，GNU 调试器，用于断点、栈帧、变量、反汇编和崩溃现场分析。
  \item[C++20] C++ 标准版本名。本册默认使用的语言和标准库能力边界。
  \item[SIMD] Single Instruction, Multiple Data，一条指令处理多个数据 lane。
  \item[AVX/AVX2] x86 向量指令扩展，常用于 CPU 后端热循环。
  \item[AVX-512] x86 512 位向量扩展，提供更宽寄存器和 mask 机制，但要验证频率、tail 和平台覆盖。
  \item[FMA] Fused Multiply-Add，一条指令完成乘加。
  \item[VNNI] Vector Neural Network Instructions，面向低精度 dot-product 的向量指令能力。
  \item[AMX] Advanced Matrix Extensions，x86 tile 级矩阵扩展，适合足够大的矩阵块计算。
  \item[PMU] Performance Monitoring Unit，硬件性能计数器来源。
  \item[IPC] Instructions Per Cycle，每周期执行指令数，需要结合内存、分支和前后端瓶颈解释。
  \item[TLB] Translation Lookaside Buffer，地址翻译缓存。
  \item[L2/L3] CPU 缓存层级。L2 通常比 L1 大且更慢，L3 通常更大并可能被多个核心共享。
  \item[LLC] Last-Level Cache，最后一级缓存。
  \item[DRAM] Dynamic Random-Access Memory，主存。
  \item[HBM] High Bandwidth Memory，高带宽显存，常见于数据中心 GPU。
  \item[NUMA] Non-Uniform Memory Access，非统一内存访问，不同节点访问成本不同。
  \item[JSON] JavaScript Object Notation，结构化文本格式，常用于 config、tokenizer、manifest 和报告。
  \item[CSV] Comma-Separated Values，逗号分隔文本表格，常用于 benchmark 和 trace 摘要。
  \item[CI] Continuous Integration，持续集成，用自动测试和构建保护回归。
  \item[I/O] Input/Output，输入输出，模型加载、日志、trace 和文件映射都会涉及。
  \item[Unicode] 统一字符集合，tokenizer 和文本规范化必须明确它和字节编码的关系。
  \item[UTF-8] Unicode 的常见变长字节编码。prompt、stop string 和 golden case 要保存原始字节边界。
  \item[KB/MB/GB] 存储容量单位。模型文件、权重、KV cache 和 arena 报告必须统一单位口径。
  \item[KiB/MiB/GiB] 二进制容量单位，分别按 1024 进位。报告内存、显存和 mmap 区间时，混用 GB 与 GiB 会造成可见误差。
  \item[mmap] 把文件或匿名内存映射到进程虚拟地址空间的系统接口。模型加载常用它减少显式读拷贝，但仍要处理页故障、生命周期和文件截断风险。
  \item[perf] Linux 性能分析工具，用采样和 PMU 事件观察热点、cache miss、分支、调度和调用栈。
  \item[objdump] 目标文件和可执行文件查看/反汇编工具，用来确认后端是否生成预期 SIMD、VNNI 或 AMX 指令。
\item[BM/BN/BK] GEMM tile 记号。BM 是 batch 行块，BN 是输出通道块，BK 是归约维度块。
\item[TILE\_N/TILE\_K] 后端 schedule 常量，分别表示 N 方向列块和 K 方向归约块大小。
  \item[量化] 用较低位宽的整数近似表示实数值，并用 scale 和 zero point 记录映射关系。
  \item[校准] 为量化选择数值范围和 scale 的过程。权重可以离线扫描，activation 和 KV cache 通常需要代表性输入或运行时策略。
  \item[Requantize] 把较宽累加结果重新映射到目标低位宽输出的步骤，通常包含固定点乘法、舍入、加 zero point 和 clamp。
  \item[Group Quantization] 按一组连续元素共享 scale 的量化方式，常用于 int4/int8 大模型权重；它降低误差但要求 kernel 正确读取 group scale。
  \item[QuantDescriptor] 描述量化格式的结构化元数据，包括 bit width、scale dtype、zero point policy、group size、axis、packed order 和 tail policy。
  \item[CalibrationRecord] 记录量化校准来源和结果的工程对象，包括校准集、算法、范围、饱和率、scale 摘要和工具版本。
  \item[Reference Converter] 用清晰标量实现把低比特字节、scale 和 zero point 还原为参考浮点值的组件，用来验证 packed order、tail policy 和 metadata 解释。
  \item[Reference] 用来定义正确结果的清晰实现，不以速度为目标。
  \item[Oracle] 判定实现输出是否满足 reference 和误差合同的比较器。
  \item[Kernel] 某个算子在特定 dtype、layout 和目标机器上的具体实现。
  \item[推理图] 算子节点和张量边组成的有向计算结构。
  \item[AI 编译器] 把模型图、张量合同、量化元数据和后端能力转换成可执行计划的编译流水线。
  \item[编译前端] 把外部输入转换成已检查内部表示的阶段。在 AI 推理引擎里，它读取模型文件、tokenizer、权重索引和量化 metadata，输出 manifest 与 typed graph。
  \item[IR] 中间表示。AI 编译器中常分为表达算子依赖的 graph IR、表达 shape/layout/dtype 的 tensor IR，以及表达 tile/vector/thread 的 loop IR。
  \item[SSA] Static Single Assignment，静态单赋值形式，每个虚拟值只定义一次，便于数据流和优化分析。
  \item[Pass] 读取一种 IR 并产出改写后 IR 或分析事实的编译步骤，例如融合 pass、layout pass、liveness pass 和 lowering pass。
  \item[Analysis Pass] 不改写 IR、只收集事实的 pass，例如 use-def、liveness、alias、shape、quant 和 cost analysis。
  \item[Transform Pass] 根据已有分析事实改写 IR 的 pass，例如融合、消除冗余 reshape、量化 lowering 和 layout 转换。
  \item[Lowering] 把较高层表示逐步降到更接近目标后端的表示，例如从算子图降到 CPU SIMD kernel 或 GPU device kernel 调用。
  \item[后端] 把已验证的 IR 或 executable plan 落到具体硬件和运行环境的实现层，例如 CPU SIMD/多线程后端或 GPU stream/kernel 后端。
  \item[Executable Plan] 编译器输出给 runtime 的可执行计划，包含 segment、kernel 选择、arena offset、外部状态绑定点和后端同步要求。
  \item[内存规划] 根据张量生命周期、大小、对齐和后端位置安排 arena、workspace、packed weight 和 KV cache 的过程。
  \item[Manifest] 模型导入后形成的机器可读清单，记录 config、tokenizer、张量位置、shape、dtype、layout 和量化元数据。
  \item[Verifier] 检查 IR、shape、dtype、量化和状态边是否满足合同的组件；失败时应给出可定位诊断，而不是让后端继续执行。
  \item[Schema Resolution] 把不同模型格式或命名方言映射到引擎内部 canonical role 的过程。
  \item[Activation Arena] 根据中间张量生命周期预先规划的一块或多块临时内存，用 offset 复用短生命周期 activation，避免热路径反复分配。
  \item[Packed Weight] 后端准备阶段从逻辑权重转换出的物理布局，通常包含 tile 重排、对齐、padding、量化 scale 布局和 kernel 读取合同。
  \item[KV Cache Quantization] 对推理过程中动态增长的 K/V 状态做低比特存储的策略。它不同于静态权重量化，必须同时考虑写入成本、metadata、长上下文误差和 decode attention 热路径。
  \item[KV cache] Key/Value cache，Transformer 解码时保存历史 token 的 K/V 张量，避免每一步重算全部历史。
  \item[QKV] Query、Key、Value。attention 用 Q 查询历史 K，再用权重加权 V。
  \item[QH/KVH/KVA] Query head 数和 key/value head 数的形状变量。GQA/MQA 下二者关系必须被 verifier 检查。
  \item[BN] Batch 或 block 相关维度记号，具体含义要看所在公式或 fixture。
  \item[MLP] Multi-Layer Perceptron，多层感知机。Transformer 中通常指 attention 后的前馈网络，由升维、激活/门控和降维投影组成。
  \item[GEMM] General Matrix-Matrix Multiplication，矩阵乘矩阵。prefill 阶段常更接近 GEMM。
  \item[GEMV] General Matrix-Vector Multiplication，矩阵乘向量。decode 阶段很多投影更接近 GEMV。
  \item[FLOP/FLOPS/TFLOP] FLOP 是一次浮点运算；FLOPS 是每秒浮点运算次数；TFLOP 表示每秒万亿次浮点运算。指标必须和 dtype、batch、内存带宽一起解释。
  \item[MFLOP/GFLOP] 百万/十亿级浮点运算量或吞吐单位，报告中要区分总量和每秒速率。
  \item[FP16/BF16/TF32/INT8] 常见低精度数据类型。FP16 是 16 位浮点；BF16 保留更多指数位；TF32 是 NVIDIA Tensor Core 常见计算格式；INT8 是 8 位整数推理常见格式。
  \item[MAC] Multiply-Accumulate，乘加次数。很多模型账本会先数 MAC，再按约定换算 FLOP。
  \item[RMSNorm] Root Mean Square Layer Normalization，用均方根缩放 hidden state 的归一化方式。
  \item[RoPE] Rotary Position Embedding，旋转位置编码，把位置信息注入 Q/K。
  \item[SwiGLU] Swish-Gated Linear Unit，一种门控前馈网络结构。
  \item[MHA] Multi-Head Attention，多头注意力，每个 query head 有独立 K/V head。
  \item[GQA] Grouped-Query Attention，多个 query head 共享一组 KV head，降低 KV cache 体积。
  \item[MQA] Multi-Query Attention，所有 query head 共享一组 KV head，进一步降低 KV cache。
  \item[QK] Attention 中 query 与 key 的乘积，用来得到注意力分数。
  \item[PV] Attention 中概率权重与 value 的乘积，用来得到上下文表示。
  \item[FlashAttention] 分块 attention 思路，通过在线 softmax 和 tile 复用减少中间矩阵写回。
  \item[Segment] executable plan 中比单个算子更粗的调度单元，绑定 arena offset、外部状态和一段 kernel 序列，用来降低运行时解释开销。
  \item[Time to First Token] 从请求进入到第一个生成 token 可用之间的延迟，主要受 tokenizer、prefill、KV cache 初始化和首轮 logits 影响。
  \item[TTFT] Time To First Token，首 token 延迟。
  \item[TPOT] Time Per Output Token，生成阶段每个输出 token 的平均时间。
  \item[p50/p95/p99] 延迟或资源占用分位数。p50 描述典型请求，p95/p99 描述尾部请求，模型服务验收必须同时报告。
  \item[SLO] Service Level Objective，服务等级目标，例如 TTFT、TPOT 或 p95 延迟阈值。
  \item[FIFO] First In First Out，先进先出。服务调度中 FIFO 简单但可能让长请求阻塞短请求。
  \item[FAQ] Frequently Asked Questions，常见问题。RAG chunking 要避免把 FAQ 的问题和答案拆散。
  \item[DDIA] Designing Data-Intensive Applications，常被用作数据密集系统设计材料简称。本册提到它时，关注 artifact、调度、日志和可恢复系统思维。
  \item[RAG] Retrieval-Augmented Generation，检索增强生成，先检索外部资料再把资料放入 prompt。
  \item[BM25] 经典稀疏关键词检索打分方法，RAG 中常和 dense vector 检索互补。
  \item[HTML] HyperText Markup Language，网页标记语言。RAG ingestion 需要把它解析成安全文本和元数据。
  \item[LLM] Large Language Model，大语言模型。
  \item[GGUF] llama.cpp 生态常见模型文件格式，打包权重、量化信息和模型元数据。
  \item[ggml] llama.cpp 生态中的张量和推理基础库。阅读它时重点看 tensor layout、量化格式、kernel dispatch 和内存规划。
  \item[llama.cpp] C/C++ 本地 LLM 推理项目，常作为 CPU/GPU 后端、GGUF、tokenizer 和量化工程的源码阅读样本。
  \item[vLLM] 高吞吐 LLM 服务系统，常用于对照 continuous batching、KV cache 管理、调度和服务 SLO。
  \item[oneDNN] CPU 深度学习算子库，提供卷积、矩阵乘、归一化等优化 primitive，可作为手写 kernel 的正确性和性能对照。
  \item[ONNX] Open Neural Network Exchange，开放神经网络交换格式，用图、算子、张量和权重描述模型。
\item[BPE] Byte Pair Encoding，一类 tokenizer 子词合并算法。
\item[BOS/EOS/PAD] BOS 是句首 token，EOS 是句尾 token，PAD 是填充 token。
\item[组合缩写] 斜杠并列的缩写组合，例如 \code{GEMM/GEMV} 或 \code{MLIR/TVM}。它提示这些概念在当前段落被放在一起比较，不代表新术语。
  \item[UNK] Unknown token，未知 token，用来表示 tokenizer 无法映射到已知词表项的输入。
  \item[OOV] Out-of-vocabulary，词表外输入。BPE 等子词 tokenizer 会降低 OOV 概率，但词表、normalization 和 byte fallback 仍要验证。
  \item[MoE] Mixture of Experts，专家混合结构。
  \item[LoRA] Low-Rank Adaptation，低秩适配，用低秩增量调整模型权重。
  \item[GPTQ] 一类训练后权重量化方法，常用近似二阶信息逐层降低量化误差。
  \item[AWQ] Activation-aware Weight Quantization，激活感知权重量化，常通过保护重要通道降低低比特权重量化损失。
  \item[NF4] NormalFloat 4-bit，一类非均匀 4-bit codebook 量化格式，通常要同时保存 codebook、block scale 和 dequant 合同。
\item[LCQI] Local CPU Quantized Inference，本书本地 CPU 量化推理实验线的项目代号。
\item[LCQI\_TOKENIZER\_V1] LCQI 教学工程里的 tokenizer 合同版本名，用来固定 vocab、special token、normalization 和 byte fallback 解释。
  \item[AST] Abstract Syntax Tree，抽象语法树，传统编译器前端常用的源码结构表示。
  \item[MLIR/TVM/XLA] 工业编译器或编译框架。本书用它们比较 IR、pass、schedule 和 lowering 边界。
  \item[CUDA/HIP] CUDA 是 NVIDIA GPU 编程接口，HIP 常用于 AMD GPU 生态。
  \item[NCCL] NVIDIA Collective Communications Library，多 GPU collective 通信库，常见于 all-reduce、broadcast 等训练或服务并行场景。
  \item[ACL] Arm Compute Library，Arm 生态常见计算库，可作为 CPU/GPU 后端算子实现来源。
  \item[PCIe] Peripheral Component Interconnect Express，主机和设备之间常见高速互连。
  \item[OOM] Out Of Memory，内存不足。
  \item[RNG] Random Number Generator，随机数生成器，采样时必须纳入可复现性边界。
  \item[NaN] Not-a-Number，浮点数中的非数字值，常提示数值计算已经失控。
  \item[KL/MSE] KL divergence 和 mean squared error，量化校准常用的分布差异或误差指标。
  \item[MMR] Maximal Marginal Relevance，最大边际相关性，检索中用于在相关性和多样性之间折中。
  \item[MRR] Mean Reciprocal Rank，平均倒数排名，用于评价第一个相关结果排得是否靠前。
  \item[PII] Personally Identifiable Information，个人可识别信息。RAG 和日志系统必须避免把它无意写入 prompt、trace 或训练样本。
  \item[UB] Undefined Behavior，C++ 未定义行为。推理后端中的越界、悬空指针、错误对齐和数据竞争都可能触发 UB。
  \item[Micro-kernel] 面向特定 dtype、layout、tile 和指令集的小型核心计算内核，例如一个 AVX2 int8 GEMV tile。它通常由外层 packing、线程调度和尾部路径包围。
  \item[RAII] C++ 中用对象生命周期管理资源的惯用法。推理后端中，文件映射、aligned buffer、device buffer、stream 和 profiler event 都应由 RAII 对象管理。
  \item[Baseline] 优化前可重复测量的基线实现。它可以不快，但必须正确、可构建、可 benchmark，并能作为后续优化的比较对象。
  \item[Tail Path] 处理维度不能整除 SIMD 宽度、tile 大小或 group size 时的边界路径。它必须被 descriptor、verifier、kernel 和测试共同覆盖。
  \item[ThreadPlan] CPU 后端对任务划分、最小任务粒度、barrier、scratch、输出写策略和亲和策略的描述。
  \item[DeviceBuffer] GPU 后端中拥有显存生命周期的 RAII 对象，避免裸 device pointer 在计划和运行时之间泄漏所有权。
  \item[SM/CU] GPU 上的基本执行资源组。NVIDIA 常称 SM，AMD 常称 CU；它们包含调度、寄存器、执行单元、shared memory 和矩阵计算单元等资源。
  \item[HBM/GDDR] GPU 常见显存类型。HBM 带宽高、封装近；GDDR 更常见于消费级 GPU。
  \item[Warp/Wavefront] GPU 中一组 lockstep 执行的线程。它类似放大的 SIMD lane 组，喜欢规则分支和连续访存。
  \item[Shared Memory] GPU SM/CU 内部的快速小容量内存，常用于 tile 复用；容量和使用量会影响 occupancy。
  \item[Tensor Core] GPU 中面向矩阵块乘加的专用计算单元，通常对 dtype、shape、layout 和对齐有要求。
  \item[Occupancy] 一个 SM/CU 上可同时驻留的 warp/block 相对上限的比例。它影响隐藏延迟的能力，但不能单独代表性能。
  \item[Stream] GPU 后端的异步执行队列。kernel launch、拷贝和 event 通常挂在 stream 上，runtime 必须用它表达同步边界。
  \item[StagingBuffer] Host 与 device 之间传输 logits、token id 或小 metadata 的中转缓冲。它的大小、同步和复用会影响 decode latency。
  \item[Performance Gate] 性能回归门禁，记录硬件、输入、baseline、阈值和指标，用来判断一次改动是否让关键路径退化。
  \item[ProfilerEvent] 统一性能事件记录，通常包含 stage、segment、backend、kernel、时间、context length、内存和可选硬件计数器。
  \item[UI] User Interface，用户界面。服务和工具章节提到 UI 时，重点是状态机、取消和回调顺序。
  \item[PR] Pull Request，代码合并请求。验收章节用它表示进入审查和门禁的变更单位。
  \item[PDF] Portable Document Format，固定版式文档格式。RAG ingestion 读取 PDF 时要处理页眉页脚、断行、表格和页码噪声。
  \item[QA] Quality Assurance 或 Question Answering。工程验收章节多指质量保证；RAG/应用章节多指问答任务。
\end{description}
